\scanpage{7}
\mspage{5}

où ce foncteur est une \uline{équivalence de catégories}: il en est en tous cas ainsi si $S\longrightarrow R$ est ``\uline{formellement lisse}'', par exemple si c'est un morphisme lisse, ou si S et R sont des spectres de corps $k_0,k$, avec k une extension \uline{séparable} de $k_0$. \deleted{Dans ce dernier cas}(Quand d'ailleurs $S\to R$ est même formellement étale, alors les deux catégories: celle des cristaux et celle des objets à stratification, deviennent \deleted{d'ailleurs} équivalentes,par le foncteur ``valeur en S''\inserted{,} à $\underline F_S$ lui-même, ce qui nous redonne l'\unsure{exmple|exemple} à la noix de 1.5 ou k est de car. nulle). Sauf erreur, on a encore une équivalence de catégories si $S\to R$ est \uline{plat et localement de présentation finie}, \inserted{(du moins si R loc.noeth., et se bornant aux Modules cohérents)} mais je n'ai pas écrit \unsure{de|la} démonstration.

1.9. \uline{\unsure{UNe|Une} digression sur la notion de stratification en caractéristique nulle}.

Quand S est lisse sur R (en fait, il suffit que S soit différentiellepent lisse sur R, i.e. le morphisme diagonal $S\to S\times_RS$ une immersion régulière), et si R est de caractéristique nulle, alors une stratification d'un Module M sur S (relativement à R) est connue quand on connait la 1-\unsure{donnexion|connexion} qu'elle définit, et on peut se donner celle-ci arbitrairement, soumise à la seule condition que le ``tenseur courbure'', qui est une certaine section de $\underline\Omega^2_{S/R}\otimes\underline{\operatorname{ENd}}(M)$, soit nul. (Cela peut aussi s'exprimer en \unsure{ddsant|disant} qu'on fait opérer le faisceau $\underline{\operatorname{Dér}}_{S/R}$ des dérivations relatives de S sur R, sur M, de façon à satisfaire aux relations habituelles, y compris celle de transformer crochet en crochet). J'ignore dans quelle mesure l\deleted{a}\inserted{'}\deleted{\unsure{propriété|proposition}} \inserted{hypothèse} de lissité est nécessaire pour cet énoncé, ou si (dans le cas lisse disons) on peut formuler un énoncé analogue pour toute catégorie fibrée $\underline F$ sur $(\mathrm{Sch})_{/R}$. Mais bien entendu, l'hypothèse de caractéristique nulle faite sur R est tout à fait essentielle. Si R est de caractéristique $p>0$, l'énoncé qui remplace le précédent (et qui sauf erreur est dû à Cartier) est que dans ce cas, la donnée d'une ``connexion sans torsion'' sur M \inserted{``compatible aux puissances \unsure{p.ièmes|pièmes}''} équivaut à une ``donnée de descente'' sur M relativement à frobenius $S\longrightarrow S^{(p/R)}$. Il y a loin de là à une \deleted{\unsure{rigidifi}} stratification !

1.10. \uline{La notion de} p-\uline{cristal et ses variantes}.

Nous supposons maintenant que $\deleted{\unsure{R}}\inserted{S}$ est de caractéristique $p>0$, \deleted{\unsure{le cas de}}
\margin{6}
